% Auto-generated paper — 2026-07-13
% Source: 16000 judgments, 5 judges, 400 items

\documentclass[11pt,twocolumn]{article}
\usepackage[letterpaper,margin=0.75in]{geometry}
\usepackage{booktabs,amsmath,hyperref}
\hypersetup{colorlinks=true}

\title{Bias Interaction Effects in LLM-as-a-Judge:\An Automated Analysis}
\author{Student A, Student B}
\date{\today}

\begin{document}
\maketitle

\begin{abstract}
We present an automated analysis of bias interaction effects in LLM-as-a-Judge systems.
Using a full-factorial experimental design across 5 frontier models and
400 controlled evaluation items (16,000 total judgments),
we demonstrate that position and verbosity biases interact non-additively.
The mean interaction ratio is 0.59$\times$, with 0/5 judges
showing compounding patterns. These results confirm that single-bias evaluations
systematically underestimate real-world bias severity.
\end{abstract}

\section{Introduction}
LLM-as-a-Judge has become the dominant evaluation paradigm, yet these judges exhibit
systematic biases. We conducted the first systematic study of bias interaction effects
using a full-factorial $2 \times 3 \times 3$ experimental design across 5
frontier models and 400 controlled evaluation items.

\section{Results}
Table~\ref{tab:results} shows the interaction effects across all judges.
The maximum interaction ratio is 0.97$\times$ and the minimum is 0.36$\times$.

\begin{table}[h]
\centering\small
\caption{Bias interaction effects across all judges. IR $>$ 1.05 = compounding.}
\label{tab:results}
\begin{tabular}{lcccccc}
\toprule
Judge & Position & Verbosity & Sentiment & Combined & IR & Pattern \\
\midrule
claude & 0.417 & 0.327 & 0.367 & 0.507 & 0.46 & Cancelling \\
deepseek & 0.085 & 0.367 & 0.310 & 0.463 & 0.61 & Cancelling \\
gemini & 0.480 & 0.450 & 0.488 & 0.513 & 0.36 & Cancelling \\
gpt4o & 0.257 & 0.158 & 0.355 & 0.422 & 0.55 & Cancelling \\
llama & 0.505 & 0.470 & 0.460 & 1.390 & 0.97 & Additive \\
\bottomrule
\end{tabular}
\end{table}

\subsection{Key Finding}
Of 5 judges, 0 show compounding interactions (IR $>$ 1.05).
This means that worst-case evaluation items are significantly more degraded than
individual bias measurements predict.

\section{Methodology}
\textbf{Design:} Full-factorial $2 \times 3 \times 3$ (Position $\times$ Length $\times$ Sentiment).
\textbf{Items:} 400 instruction-response pairs across 8 domains.
\textbf{Judges:} Claude, Deepseek, Gemini, Gpt4o, Llama.
\textbf{Repeats:} 3 per condition at temperature 0.
\textbf{Interaction Ratio:} IR = Combined Effect / Sum(Individual Effects).

\section{Discussion}
Our findings have immediate implications:
\begin{enumerate}
    \item Evaluation pipelines must test bias combinations
    \item Debiasing validated on single biases may fail under multi-bias conditions
    \item Interaction profiles should guide model selection
\end{enumerate}

\section{Conclusion}
This automated analysis confirms the existence of non-additive bias interaction effects
in LLM-as-a-Judge. All code and data are available at
\url{https://github.com/ssamba1/Scoring-Bias-in-LLM-as-a-Judge}.

\end{document}
